%------------------------------------------------------------------------------%
% \chapter{Revisão}
%------------------------------------------------------------------------------%

%------------------------------------------------------------------------------%
\begin{exercicios}{Expressões Racionais}
    
    
\begin{exercicio}
  Simplifique as expressões.
  \begin{mathlist}[2]
      \item \frac{x^2 + x - 2}{x^2 - 4}
      \item \frac{2a^2 - 3ab - 9b^2}{2ab^2 + 3b^3}
      \item \frac{12t^2 + 12t + 3}{4t^2 - 1}
      \item \frac{3(4x - 1) - 4(3x + 1)}{(4x - 1)^2}
  \end{mathlist}    
  \begin{answers}
      \begin{mathlist}[2]
        \item \frac{x - 1}{x - 2}
        \item \frac{a - 3 b}{b^{2}}
        \item \frac{3 \left(2 t + 1\right)}{2 t - 1}
        \item - \frac{7}{\left(4 x - 1\right)^{2}}
    \end{mathlist}
  \end{answers}
\end{exercicio}

\begin{exercicio}
  Efetue as operações indicadas e simplifique cada expressão.
  \begin{mathlist}
      \item \frac{2a^2 - 2b^2}{b - a} \, \frac{4a + 4b}{a^2 + 2ab + b^2}
      \item \frac{x^2 - 6x + 9}{x^2 - x - 6} \, \frac{3x + 6}{2x^2 - 7x + 3}
      \item \frac{3x^2 + 2x - 1}{2x + 6} \div \frac{x^2 - 1}{x^2 + 2x - 3}
      \item \frac{58}{3 (3t + 2)} + \frac{1}{3}
      \item \frac{2x}{2x - 1} - \frac{3x}{2x + 5}
      \item \frac{-xe^x}{x + 1} + e^x
      \item \frac{4}{x^2 - 9} - \frac{5}{x^2 - 6x + 9}
      \item \frac{x}{1 - x} + \frac{2x + 3}{x^2 - 1}
      \item \frac{1 + \frac{1}{x}}{1 - \frac{1}{x}}
      \item \frac{4x^2}{\sqrt{2x^2 + 7}} + \sqrt{2x^2 + 7}
      \item 6 (2x + 1)^2 \sqrt{x^2 + x} + \frac{(2x +1)^4}{2 \sqrt{x^2 + x}}
      \item \frac{2x (x + 1)^{-1/2} - (x + 1)^{1/2}}{x^2}
      \item \frac{(x^2 + 1)^{1/2} - 2x^2 (x^2 + 1)^{-1/2}}{1 - x^2}
      \item \frac{(2x + 1)^{1/2} - (x + 2) (2x + 1)^{-1/2}}{2x + 1}
  \end{mathlist}        
  \begin{answers}
      \begin{mathlist}[2]
        \item -8
        \item \frac{3}{2 x - 1}
        \item \frac{3 x}{2} - \frac{1}{2}
        \item \frac{t + 20}{3 t + 2}
        \item \frac{x \left(13 - 2x\right)}{\left(2 x - 1\right) \left(2 x + 5\right)}
        \item \frac{e^{x}}{x + 1}
        \item - \frac{x + 27}{x^{3} - 3 x^{2} - 9 x + 27}
        \item \frac{- x^{2} + x + 3}{x^{2} - 1}
        \item \frac{x + 1}{x - 1}
        \item \frac{7+6\,x^2}{\sqrt{2\,x^2+7}}
        \item \frac{\left(1+2x\right)^2\left(1+16x+16x^2\right)}{2\sqrt{x\left(x+1\right)}}
        \item \frac{x - 1}{x^{2} \sqrt{x + 1}}
        \item \frac{1}{\sqrt{x^{2} + 1}}
        \item \frac{x - 1}{\left(2 x + 1\right)^{\frac{3}{2}}}
    \end{mathlist}
  \end{answers}
\end{exercicio}


    
%    \begin{exercicio}
%        Simplifique as expressões.
%        \begin{mathlist}
%            \item \frac{x^2 + x - 2}{x^2 - 4}
%            \item \frac{2a^2 - 3ab - 9b^2}{2ab^2 + 3b^3}
%            \item \frac{12t^2 + 12t + 3}{4t^2 - 1}
%            \item \frac{(4x - 1) (3) - (3x + 1) (4)}{(4x - 1)^2}
%        \end{mathlist}    
%    \end{exercicio}
%    
%    \begin{exercicio}
%        Efetue as operações indicadas e simplifique cada expressão.
%        \begin{mathlist}
%            \item \frac{2a^2 - 2b^2}{b - a} \, \frac{4a + 4b}{a^2 + 2ab + b^2}
%            \item \frac{x^2 - 6x + 9}{x^2 - x - 6} \, \frac{3x + 6}{2x^2 - 7x + 3}
%            \item \frac{3x^2 + 2x - 1}{2x + 6} \div \frac{x^2 - 1}{x^2 + 2x - 3}
%            \item \frac{58}{3 (3t + 2)} + \frac{1}{3}
%            \item \frac{2x}{2x - 1} - \frac{3x}{2x + 5}
%            \item \frac{-xe^x}{x + 1} + e^x
%            \item \frac{4}{x^2 - 9} - \frac{5}{x^2 - 6x + 9}
%            \item \frac{x}{1 - x} + \frac{2x + 3}{x^2 - 1}
%            \item \frac{1 + \frac{1}{x}}{1 - \frac{1}{x}}
%            \item \frac{4x^2}{2 \sqrt{x^2 + 7}} + \sqrt{2x^2 + 7}
%            \item 6 (2x + 1)^2 \sqrt{x^2 + x} + \frac{(2x +1)^4}{2 \sqrt{x^2 + x}}
%            \item \frac{2x (x + 1)^{-1/2} - (x + 1)^{1/2}}{x^2}
%            \item \frac{(x^2 + 1)^{1/2} - 2x^2 (x^2 + 1)^{-1/2}}{1 - x^2}
%            \item \frac{(2x + 1)^{1/2} - (x + 2) (2x + 1)^{-1/2}}{2x + 1}
%        \end{mathlist}        
%    \end{exercicio}
    
    \begin{exercicio}
        Racionalize o denominador de cada expressão.
        \begin{mathlist}[2]
            \item \frac{1}{\sqrt{3} - 1}
            \item \frac{1}{\sqrt{x} - \sqrt{y}}
            \item \frac{a}{1 - \sqrt{a}}
            \item \frac{\sqrt{a} + \sqrt{b}}{\sqrt{a} - \sqrt{b}}
        \end{mathlist}    
    \end{exercicio}
    
    \begin{exercicio}
        Racionalize o numerador de cada expressão.
        \begin{mathlist}[2]
            \item \frac{\sqrt{x}}{3}
            \item \frac{\sqrt[3]{y}}{x}
            \item \frac{1 - \sqrt{3}}{3}
            \item \frac{1 + \sqrt{x + 2}}{\sqrt{x + 2}}
        \end{mathlist}        
    \end{exercicio}
    
    \begin{exercicio}
        Determine se as afirmações são verdadeiras ou falsas. Se verdadeiras, explique porquê. Se falsas, justifique através de um exemplo.
        \begin{alphalist}
            \item Se $b^2 - 4ac > 0$, então $ax^2 + bx + c = 0$, $a \neq 0$, tem duas raízes reais.
            \item Se $b^2 - 4ac < 0$, então $ax^2 + bx + c = 0$, $a \neq 0$, não tem raízes reais.
            \item \( \frac{a}{b + c} \left(\frac{a}{b} + \frac{a}{c}\right)\)
            \item \( \sqrt{(a + b) (b - a)} = \sqrt{b^2 - a^2}\) para todos os números reais $a$ e $b$.
        \end{alphalist}                                        
    \end{exercicio}
    
    \begin{exercicio}
        Encontre os valores de $x$ que satisfazem a(s) desigualdade(s).
        \begin{mathlist}
            \item - x + 3 \leq 2x + 9
            \item x - 3 > 2 \qquad \text{ou} \qquad x + 3 < -1
            \item 2x^2 \, > \, 50
        \end{mathlist}
    \end{exercicio}
    
    \begin{exercicio}
        Calcule as expressões dadas.
        \begin{mathlist}
            \item \abs{-5 + 7} + \abs{-2}
            \item \abs{2\pi - 6} - \pi
            \item \abs{\sqrt{3} - 4} + \abs{4 - 2\sqrt{3}}
        \end{mathlist}    
        \begin{answers}
            \begin{mathlist}[3]
                \item 4
                \item \pi -6
                \item 8 - 3 \sqrt{3}
            \end{mathlist}
        \end{answers}
    \end{exercicio}
    
    \begin{exercicio}
        Calcule as expressões dadas.
        \begin{mathlist}[2]
            \item \left(\frac{9}{4}\right)^{3/2}
            \item \frac{5^6}{5^4}
            \item \frac{(3 \times 2^{-3}) (4 \times 3^5)}{2 \times 9^3}
            \item \frac{3 \sqrt[3]{54}}{\sqrt[3]{18}}
        \end{mathlist}        
        \begin{answers}
            \begin{mathlist}[4]
                \item \sfrac{27}{8}
                \item 25
                \item \sfrac{1}{4}
                \item 3 \sqrt[3]{3}
            \end{mathlist}
        \end{answers}
    \end{exercicio}
    
    \begin{exercicio}
        Simplifique as expressões dadas.
        \begin{mathlist}[2]
            \item \dfrac{4(x^2 + y)^3}{x^2 + y}
            \item \dfrac{a^6 B^{-5}}{(a^3b^{-2})^{-3}}
            \item \dfrac{\sqrt[4]{16x^5yz}}{\sqrt[4]{81xyz^5}} \) \\[2mm] $x$, $y$ e $z$ positivos
            \item \left(\dfrac{3xy^2}{4x^3y}\right)^{-2} \left(\dfrac{3xy^3}{2x^2}\right)^3    
        \end{mathlist}  
        \begin{answers}
            \begin{mathlist}[4]
                \item \null\!\! 4\! \left(\! x^{2}\! +\! y \!\right)^{2}
                \item \frac{a^{15}}{B^{5} b^{6}}
                \item \frac{2x}{3z}
                \item 6 x y^{7}
              \end{mathlist}
          \end{answers}          
    \end{exercicio}
    
    \begin{exercicio}
        Fatore as expressões dadas.
        \begin{mathlist}[2]
            \item -2 \pi^2 r^3 + 100 \pi r^2
            \item 2v^3w + 22 vw^3 + 2u^2 vw
            \item 12t^3 - 6t^2 - 18t
        \end{mathlist}        
        \begin{answers}
            \begin{mathlist}[2]
                \item 2 \pi r^{2} \left(- \pi r + 50\right)
                \item 2 v w \left(u^{2} + v^{2} + 11 w^{2}\right)
                \item 6 t \left(2 t^{2} - t - 3\right)
            \end{mathlist}
        \end{answers}
    \end{exercicio}
    
    \begin{exercicio}
        Faça as operações indicadas e simplifique as expressões dadas.    
        \begin{mathlist}
            \item \frac{6x}{2(3x^2 + 2)} + \frac{1}{4(x + 2)}
            \item \frac{2}{3} \left(\frac{4x}{2x^2 - 1}\right) + 3 \left(\frac{3}{3x - 1}\right)
            \item \frac{-2x}{\sqrt{x + 1}} + 4\sqrt{x + 1}
        \end{mathlist}    
        \begin{answers}
            \begin{mathlist}[2]
                \item \frac{15 x^{2} + 24 x + 2}{4 \left(x + 2\right) \left(3 x^{2} + 2\right)}
                \item \frac{78 x^{2} - 8 x - 27}{3 \left(3 x - 1\right) \left(2 x^{2} - 1\right)}
                \item \frac{2 \left(x + 2\right)}{\sqrt{x + 1}}
            \end{mathlist}
        \end{answers}
    \end{exercicio}
    
    \begin{exercicio}
        Racionalize o denominador.
        \[\frac{\sqrt{x} - 1}{2 \sqrt{x}}\]        
    \end{exercicio}
    
\end{exercicios}

%------------------------------------------------------------------------------%
